\documentclass[10pt,a4paper]{amsart}
\usepackage{amsmath,amssymb,amsthm,graphicx}
\newcommand{\quand}{\quad\text{and}\quad}
\newcommand{\irr}{\mathrm{irr}}
\newcommand{\red}{\mathrm{red}}
\newcommand{\sym}{\mathrm{sym}}

\newenvironment{smallpmatrix}{\big(\begin{smallmatrix}}{\end{smallmatrix}\big)}

\newcommand\mto{\mathchoice{\longmapsto}{\mapsto}{\mapsto}{\mapsto}}
\newcommand{\psfrac}[2]{\sfrac{(#1)}{#2}}
\newcommand{\spfrac}[2]{\sfrac{#1}{(#2)}}

\newtheorem{theo}{Theorem}[section]
\newtheorem{theorem}[theo]{Theorem}
\newtheorem{coro}[theo]{Corollary}
\newtheorem{lemma}[theo]{Lemma}
\newtheorem{prop}[theo]{Proposition}
\newtheorem{conj}[theo]{Conjecture}

\theoremstyle{definition}
\newtheorem{defi}[theo]{Definition}
\newtheorem{question}[theo]{Question}
\newtheorem{example}[theo]{Example}
\newtheorem{rmk}[theo]{Remark}
\newtheorem*{rem}{Remark}

\newcommand{\proj}{\mathbf{P}}

\newcommand{\R}{\mathbb{R}}
\newcommand{\C}{\mathbb{C}}
\newcommand{\N}{\mathbb{N}}
\newcommand{\Hyp}{\mathbb{H}}

\newcommand{\cC}{\mathcal{C}}

\newcommand{\Cone}{\mathrm{Cone}}
\newcommand{\diag}{\mathrm{diag}}
\newcommand{\SL}{\mathrm{SL}}
\newcommand{\SO}{\mathrm{SO}}

\newcommand{\fa}{\mathfrak{a}}
\newcommand{\frakg}{\mathfrak{g}}
\newcommand{\fp}{\mathfrak{p}}
\newcommand{\fk}{\mathfrak{k}}

\newcommand{\norm}[1]{\Vert #1 \Vert}
\newcommand{\scal}[2]{\langle #1\, |\, #2 \rangle }

\DeclareMathOperator{\Card}{Card}
\DeclareMathOperator{\DimH}{DimH}
\DeclareMathOperator{\Sym}{Sym}
\DeclareMathOperator{\Span}{Span}
\DeclareMathOperator{\ad}{ad}

\let\epsilon\varepsilon
\let\op\iota
\let\phi\varphi
\let\p\varphi
\let\G\Gamma
\let\g\gamma

\let\Log\log

\title{Hausdorff dimension n-2 of the symmetric boundary of a strictly convex projective domain with C1 boundary}
\author{Olivier Glorieux, Daniel Monclair, Nicolas Tholozan}
\date{doi:10.5802/jep.241; CC BY 4.0}
\begin{document}
\maketitle
\setcounter{section}{2}
\setcounter{theo}{54}
\noindent\textit{Editorial note.} All delivered mathematical text is the authors' own, copied from the fixed publisher source (Journal de l'Ecole polytechnique - Mathematiques 10 (2023), DOI 10.5802/jep.241, statement source lines 1008-1012 and proof source lines 1016-1061); no statement, proof step or formula was rewritten, shortened or removed.\par\medskip
\noindent\textit{Editorial note.} The publisher class cedram is not available here and is replaced by a plain amsart wrapper; the preamble consists of the authors' own theorem declarations and notation macros copied verbatim from source lines 35-45 and 67-117. The journal's title-page styling block, which loads a journal logo file, is deliberately not retained because it belongs to the publisher layout and not to the mathematics; no mathematical macro was invented, renamed or added.\par\medskip
\noindent\textit{Editorial note.} The excerpt carries no section heading, so the section counter is fixed at 2 and the authors' shared theorem counter at 54 before the body; the counted statement therefore prints as Theorem 2.55, which is its number in the article under the authors' own declaration newtheorem\{theo\}\{Theorem\}[section]. The proof environment prints no number of its own, and the two context paragraphs carry none either.\par\medskip
\noindent\textit{Editorial note.} The proof environment is not adjacent to the statement: the authors' remark at source line 1014 sits between them, and the proof environment itself names the statement it proves, 'Proof of Theorem \textbackslash{}ref\{theo - HD Symmetric boundary convex\}'. That remark is carried as context, and the pairing is the authors' own explicit designation.\par\medskip
\noindent\textit{Editorial note.} The notation x* for the tangent hyperplane at a point of the boundary, used by the statement, is introduced in the context paragraph at source lines 1003-1005 and again inside the statement's own display; both are authors' text.\par\medskip
\noindent\textit{Editorial note.} No figure, table or bibliography entry is referenced by the retained spans, so none is carried.\par\medskip
\noindent\textit{Allowed standard prerequisites (editorial).}\begin{itemize}
\item Hausdorff dimension and Hausdorff measure of a subset of a Euclidean space. DimH(A) is the infimum of the s for which the s-dimensional Hausdorff measure of A vanishes; a bi-Lipschitz image of A has the same dimension, and a graph of a Lipschitz function on an open subset of R\textasciicircum{}m has dimension m.. the authors' own definition and conventions at source lines 309-331. It is the invariant being computed and the tool used at the end of the proof, where the set is shown to be a 1-Lipschitz graph over a spacelike subspace of dimension n-2.
\item Duality and affine charts in real projective space. Every point x of a convex domain boundary has a tangent hyperplane x*, the dual projective space parametrizes hyperplanes, and after choosing two boundary points x\_0, x\_infinity with covectors alpha\_0, alpha\_infinity the affine patch P(R\textasciicircum{}n) minus x\_infinity\textasciicircum{}* is identified with V x R, V = x\_0\textasciicircum{}* cap x\_infinity\textasciicircum{}*, by (v,t) -> [v + t v\_infinity + v\_0].. the authors' own setup at source lines 1017-1023. Only this elementary projective duality and the two explicit affine charts are assumed; the formula for the dual point and the dimension bound are retained obligations.
\item Strict convexity of a C\textasciicircum{}1 function and Brouwer's invariance of domain. For a strictly convex C\textasciicircum{}1 function f one has d\_u f(v-u) < f(v) - f(u) for u different from v; an injective continuous map from an open subset of R\textasciicircum{}m into R\textasciicircum{}m is open.. used at source lines 1055-1060. Strict convexity supplies the strict positivity of the quadratic form on difference vectors, and invariance of domain converts injectivity of the orthogonal projection onto a spacelike subspace into the graph representation that yields the 1-Lipschitz bound.
\end{itemize}\medskip
A particular case of strongly projectively convex-cocompact group in $\SL(n,\R)$ is a group acting properly discontinuously and cocompactly on a Gromov-hyperbolic convex set. Crampon proved that the Hilbert critical exponent of such a group is at most $n-2$. we prove here that the Hausdorff dimension of $\Lambda_\Gamma^{\sym}$ equals $n-2$, so that our main theorem does recover Crampon's result.

If $\Gamma$ acts properly discontinuously and cocompactly on a Gromov hyperbolic convex domain $\Omega$, then we have $\Lambda_\Gamma = \partial \Omega$, $\Lambda_\Gamma^* = \partial \Omega^*$ (the set of hyperplanes tangent to $\partial\Omega$) and $\Lambda^{\sym}_\Gamma \subset \Lambda_\Gamma \times \Lambda_\Gamma^*$ is the graph of the homeomorphism mapping $x\in \partial\Omega$ to $x^* := T_x \partial \Omega$.

\begin{theorem} \label{theo - HD Symmetric boundary convex}
Let $\Omega$ be a strictly convex domain of $\proj(\R^n)$ with $\mathcal C^1$ boundary. Then
\[\partial \Omega^{\sym} := \{(x,x^*)\mid x\in \partial \Omega\} \subset \partial \Omega \times \partial \Omega^*\]
has Hausdorff dimension $n-2$.
\end{theorem}

Note that, while $\partial \Omega$ and $\partial\Omega^*$ are $\mathcal C^1$ spheres of dimension $n-2$, the homeomorphism $x\mto x^*$ is not $\mathcal C^1$ unless $\partial \Omega$ is $\mathcal C^2$ (which is very constraining for divisible convex sets). It is thus not immediate that its graph has Hausdorff dimension $n-2$ in general.

\begin{proof}[Proof of Theorem \ref{theo - HD Symmetric boundary convex}]
Fix two distinct points $x_0,x_\infty\in \partial \Omega$, and consider vectors $v_0\!\in\!x_0$, $v_\infty\!\in\!x_\infty$ and linear forms $\alpha_0\!\in\!x_0^*$, $\alpha_\infty\!\in\!x_\infty^*$ such that $\alpha_0(v_\infty)=\alpha_\infty(v_0)=\nobreak1$. Let $V=x_0^*\cap x_\infty^*\subset \R^n$, so that the affine patch $\proj(\R^n)\setminus x_\infty^*$ can be identified with $V\times \R$ through the map $(v,t)\mto [v+tv_\infty+v_0]$.

Since $\Omega$ is strictly convex and has $\mathcal C^1$ boundary, under this identification $\partial\Omega\setminus \{x_\infty\}$ is the graph of a $\mathcal C^1$ strictly convex function $f:V\to \R$.

Let $p_{0,\infty}:\R^n\to V$ be the projection with kernel $x_0\oplus x_\infty$, so that we also get an identification of the affine patch $\proj\left((\R^n)^*\right)\setminus x_0^\perp$ with $V^*\times \R$ through the map $(\alpha,t)\mto [\alpha\circ p_{0,\infty} + t \alpha_\infty + \alpha_0]$.

Under this identification, the dual point of $(u,f(u))\in \partial \Omega\setminus \{x_\infty\}$ is
\[
(u,f(u))^*=(-d_uf,d_uf(u)-f(u))\in \partial\Omega^*\setminus \{x_\infty^*\}.
\]
It follows that the map
\begin{align*}
\partial\Omega^*\setminus \{x_\infty^*\} & \to V^* \\
(u,f(u))^* & \mto -d_uf
\end{align*}
is just the projection on the first coordinate of a $\mathcal C^1$ submanifold of $V^*\times \R$, so it is~$\mathcal C^1$ (even though it is the composition of the non $\mathcal C^1$ maps $x^*\mto x$ and $v\mto d_vf$).

The projection
\begin{align*}
\pi: \left(\partial \Omega \setminus \{x_\infty\}\right)\times \left(\partial \Omega^*\setminus \{x_\infty^*\}\right) & \to V\times V^*\\
\left((u, f(u)), (v,f(v))^*\right) & \mto (u, \mathrm d_v f)
\end{align*}
is a $\mathcal C^1$ isomorphism. It is thus enough to prove that
\[\pi(\partial \Omega^{\sym} \setminus \{(y,y^*)\}) = \{(u, \mathrm d_u f)\mid u\in V\}\]
has Hausdorff dimension $d-2$.

Now let us equip $V \times V^*$ with the symmetric bilinear form
\[\langle (u,\alpha), (v,\beta)\rangle = \alpha(v) + \beta(u).\]
This form is non-degenerate of signature $(n-2,n-2)$, and we claim that
\[
\pi(\partial \Omega^{\sym} \setminus \{(y,y^*)\})
\]
is a spacelike topological submanifold of dimension $n-2$, hence a Lipschitz manifold.

Indeed, for every $u\neq v \in V$, we have
\begin{align*}
\bigl\langle (v-u, \mathrm d_v f - \mathrm d_u f), (v-u, \mathrm d_v f - \mathrm d_u f)\bigr\rangle &= -2(\mathrm d_u f(v-u) + \mathrm d_v f(u-v)).
\end{align*}
By strict convexity of $f$, we have
\[\mathrm d_u f(v-u) < f(v) - f(u)\quand \mathrm d_v f(u-v) < f(u) - f(v),\]
from which we get
\[\bigl\langle (v-u, \mathrm d_v f - \mathrm d_u f), (v-u, \mathrm d_v f - \mathrm d_u f)\bigr\rangle > 0.\]

Now, let $W$ be a spacelike subspace of $V\times V^*$ of dimension $n-2$. The previous computation implies that the orthogonal projection of $\pi(\partial \Omega^{\sym} \setminus \{(x_\infty, x_\infty^*)\})$ to $W$ is injective. Since $\pi(\partial \Omega^{\sym} \setminus \{(x_\infty, x_\infty^*)\})$ is a $\mathcal C^1$ manifold of dimension $n-2$, this projection is open (by Brouwer's invariance of domain), and $\pi(\partial \Omega^{\sym} \setminus \{(x_\infty, x_\infty^*)\})$ is the graph of a function from an open domain of $W$ to $W^\perp$. Finally, since $\pi(\partial \Omega^{\sym} \setminus\nobreak \{(x_\infty, x_\infty^*)\})$ is spacelike, this graph is $1$-Lipschitz (when $W$ is endowed with $\langle \cdot , \cdot \rangle$ and $W^\perp$ with $-\langle \cdot , \cdot \rangle$). We thus conclude that $\pi(\partial \Omega^{\sym} \setminus \{(x_\infty, x_\infty^*)\})$ is a Lipschitz manifold of dimension $n-2$, hence $\partial \Omega^{\sym}$ has Hausdorff dimension $n-2$.
\end{proof}

\end{document}
